\chapter{Development}

\section{Build targets}

The kernel builds for three architectures via custom target JSON specs in
\texttt{target\_specs/}:

\begin{itemize}
    \item \texttt{aarch64-unknown-none-catten.json}
    \item \texttt{x86\_64-unknown-none-catten.json}
    \item \texttt{riscv64gc-unknown-none-catten.json} (planned, not actively
        maintained)
\end{itemize}

Build command (headless, no display dependency):

\begin{verbatim}
cargo build --package catten \
  --target target_specs/aarch64-unknown-none-catten.json \
  --no-default-features --features acpi
\end{verbatim}

The \texttt{display} feature requires a C library for the flanterm framebuffer
terminal; headless boot uses serial output only (AArch64: PL011 UART; x86-64:
16550 COM1). AArch64 boots via \texttt{scripts/boot-smp1.sh} (single-core) or
\texttt{scripts/boot-aarch64.sh}. x86-64 boots via
\texttt{scripts/boot-x86.sh} (requires \texttt{-cpu max} for
\texttt{RDTSCP}/\texttt{FSGSBASE}).

\section{Self-tests}

There is no host-side \texttt{cargo test}; all tests are \textbf{boot-time
kernel self-tests} ($\sim$whitebox integration tests) run from
\texttt{self\_test::run\_self\_tests()} in \texttt{main.rs:104}, called after
\texttt{INIT\_BARRIER.wait()}. Each test uses \texttt{assert!()} /
\texttt{assert\_eq!()} (panic = failure). Current test suite (both arches
pass):

\begin{itemize}
    \item \texttt{self\_test/completion.rs} --- buffer ownership, cancel /
        deferred-reclaim, backpressure
    \item \texttt{self\_test/syscall.rs} --- dispatch table routing via
        synthetic TrapFrame
    \item \texttt{self\_test/ipi.rs} --- bounded queue, backpressure,
        closure dispatch
    \item \texttt{self\_test/shard.rs} --- \texttt{ShardLocal<T>} owner-check
        / re-entrancy guard, \texttt{ShardMailbox<M>} send/recv
    \item \texttt{self\_test/el0.rs} --- real-EL0 user thread (AArch64 only)
    \item \texttt{self\_test/cq.rs} --- CQ ring write/drain/overflow
    \item \texttt{self\_test/cq\_completion.rs} --- CQ ring integrated with
        completion subsystem
\end{itemize}

\section{Async-syscall demo}

\texttt{crates/catten/src/demo.rs} spawns a coordinator + worker thread from
\texttt{bsp\_main} (after the scheduler is active). The coordinator calls
\texttt{submit\_worker}, which spawns a worker and returns a capability that
fires when the worker exits; the coordinator blocks on \texttt{wait(cap)}.
The worker sleeps, returns, and the exit-observer completes the cap, waking
the coordinator. This exercises the full submit--async-work--complete--wake
loop.

Output at boot:
\begin{verbatim}
Testing Complete. All Tests Passed!
[async-demo] coordinator: submitted worker, now awaiting completion...
[async-demo] worker: performing async work (sleep 50ms)
[async-demo] worker: work finished, exiting (completion fires on exit)
[async-demo] coordinator: COMPLETION RECEIVED, result Ok(42)
[async-demo] SUCCESS: async syscall round-trip complete (via thread-exit 
             observer)
\end{verbatim}

\section{CI}

GitHub Actions (\texttt{.github/workflows/ci.yml}) builds both architectures
with \texttt{-D warnings} and runs an automated QEMU boot gate on the
\texttt{dev} branch (grepping for ``Testing Complete. All Tests Passed!'').

\section{Where to start reading}

\begin{itemize}
    \item \texttt{crates/catten/src/completion/mod.rs} --- the completion
        subsystem (capability table, submit/complete/poll/wait/cancel/close,
        \texttt{submit\_worker}, exit-observer).
    \item \texttt{crates/catten/src/completion/cq.rs} --- the CQ ring
        buffer.
    \item \texttt{crates/catten/src/syscall/mod.rs} --- AArch64 syscall
        dispatch table.
    \item \texttt{crates/catten/src/demo.rs} --- the end-to-end async demo.
    \item \texttt{crates/catten/src/cpu/scheduler/} --- threads, wakers,
        \texttt{RoundRobin}, \texttt{block\_thread}.
    \item \texttt{crates/catten/src/cpu/multiprocessor/ipi.rs} --- bounded
        IPI queues and the \texttt{IpiRpc} enum.
    \item \texttt{crates/catten/src/cpu/scheduler/threads/mod.rs:63--69} ---
        the completion-capability design comment.
\end{itemize}

\endinput
